\documentclass[11pt,a4paper]{article}
\usepackage[utf8]{inputenc}\usepackage[T1]{fontenc}\usepackage[italian]{babel}
\usepackage{amsmath,amssymb,mathrsfs,booktabs,geometry,hyperref}\geometry{margin=2.5cm}
\title{UMCH: dalla curvatura worldline a una ipotesi operativa multiscala dello spaziotempo\\Sintesi critica, risultati negativi e gate di identificabilità}
\author{Pepe Donato\\Independent Researcher}\date{2026}
\begin{document}\maketitle
\begin{abstract}
Riorganizziamo l'Ipotesi della Curvatura Minima Universale (UMCH) dall'assunto iniziale sulla worldline fino alla formulazione operativa multiscala. Una particella geodetica può avere $a^\mu=0$ in spaziotempo curvo, perciò il precedente core pointwise è \verb|SUPERSEDED_AS_CORE| e i modelli fissi A/B sono \verb|CONTRADICTED_UNDER_ASSUMPTIONS| nel loro scope. Il nuovo core usa regioni finite, frame fisico, vettore grezzo di risposte e modello nullo Minkowski. Dimostriamo che $\ell_0$ è \verb|ELL0_STRUCTURALLY_NON_IDENTIFIABLE_UNDER_CURRENT_FAMILIES|. Una famiglia con turnover rende $\ell_0$ identificabile in principio se la shape è fissata indipendentemente, ma resta \verb|MATHEMATICAL_IDENTIFIABILITY_CANDIDATE_ONLY|: la geometria locale giustifica $p=2$ solo in canali area-normalizzati e non deriva il coefficiente esponenziale $q$. Stato complessivo: \verb|UNPROVEN| e \verb|NO_POSITIVE_DETECTION_CLAIM|.
\end{abstract}

\section{Domanda scientifica e ledger}\label{sec:u-question}
La domanda non è se soluzioni curve della relatività generale esistano, ma se lo spaziotempo realizzato mostri una risposta geometrica operativa positiva su ogni regione ammissibile oltre una scala universale. UMCH-U-0001: questa è una ipotesi falsificabile, non un risultato. Il ledger separa geometria nota, definizioni del progetto, risultati negativi, candidati matematici e osservazioni; nessun dato reale è analizzato qui.

\section{Ipotesi iniziale e correzione fisica}\label{sec:u-history}
Il programma iniziale associava una scala $\kappa_0$ a curvatura o accelerazione pointwise della worldline. Ma per moto libero
\begin{equation}\label{eq:u-geodesic}a^\mu=u^\nu\nabla_\nu u^\mu=0
\end{equation}
può valere mentre il tensore di Riemann è non nullo. UMCH-U-0002: nessuna forza di reazione è necessaria o introdotta. Il core worldline è \verb|SUPERSEDED_AS_CORE|, non cancellato.

\section{Risultati negativi storici}\label{sec:u-negative-history}
Le analisi dei modelli pointwise fissi A e B hanno escluso geodetiche o prodotto settori degeneri/non identificabili sotto le assunzioni dichiarate. UMCH-U-0003: A/B sono \verb|CONTRADICTED_UNDER_ASSUMPTIONS|; il risultato combinato è un no-go condizionale, non una confutazione di ogni UMCH. La storia resta artefatto tecnico riproducibile.

\section{Ipotesi spacetime riformulata}\label{sec:u-hypothesis}
Sia $(\mathcal M,g_{\rm real})$ lo spaziotempo realizzato, non globalmente isometrico a Minkowski. Introduciamo $\ell_0>0$. Il survey domain $D_{\rm survey}$ e le scale $\ell$ sono fissati indipendentemente da $\ell_0$ e dalle risposte; un modello supportato solo per $x\geq1$ deve soddisfare $\ell_0\leq\ell_{\min}$ su tutti i campioni, altrimenti è \verb|DOMAIN_INCONSISTENT| e nessun campione può essere rimosso post-fit. UMCH-U-0004: il core \verb|UNPROVEN| propone, per protocollo preregistrato e frame risolto, una risposta dimensionless sopra una soglia scelta prima dei dati. La regola elimina selezione circolare, ma non identifica $\ell_0$.

\section{Regioni e frame fisico}\label{sec:u-frame}
Una regione deve essere regolare per il calcolo, causalmente testabile e associata a un estimatore di scala preregistrato. Il frame primario è l'autovettore timelike futuro unico di $T^\mu{}_{\nu}$. Nel vuoto una continuazione CMB richiede anchor, campo unitario, trasporto covariante, famiglia di cammini, unicità e copertura. UMCH-U-0005: requisiti mancanti producono \verb|FRAME_UNRESOLVED|, mai conferma. Un frame risolto è necessario, non sufficiente.

\section{Campo di risposta tensoriale/operatoriale}\label{sec:u-response}
L'oggetto primario è una sezione $\mathscr R[\Omega,\ell,u_{\rm phys},\mathcal P;g]$ di un fibrato di risposte channel-native: endomorfismi tidal/magnetic, mappe di olonomia, matrici di residui clock, mappe di Jacobi null e campi di congruenza. Il registro
\begin{equation}\label{eq:u-vector}\mathbf R=(R_{\rm tidal},R_{\rm mag},R_{\rm hol},R_{\rm clock},R_{\rm null},R_{\rm cong})\end{equation}
indicizza famiglie, non sei scalari indipendenti. Le parti Weyl seguono \cite{Bonnor1995,MaartensBassett1998}. UMCH-U-0006: stato \verb|RATIFIED_PRIMARY_RESEARCH_OBJECT|; identificabilità \verb|OPERATOR_IDENTIFIABILITY_NOT_YET_PHYSICALLY_DERIVED|. Invarianti candidati includono rapporti spettrali, spectral flow, rango/nullità, anisotropia e classi/singular values di holonomy. Il controesempio collineare è \verb|PROJECTIVE_NON_IDENTIFIABLE_COLLINEAR|.

\section{Protocolli operativi}\label{sec:u-protocols}
$R_{\rm tidal}$ e $R_{\rm mag}$ usano RMS frame-normalizzati; $R_{\rm hol}$ usa circuiti finiti preregistrati; $R_{\rm clock}$ usa residui di tempo proprio dopo baseline fissata \cite{PetitWolf2005}; $R_{\rm null}$ usa Jacobi/shear/convergenza a geometria fissata nel contesto null standard \cite{Sachs1961}; $R_{\rm cong}$ usa espansione/shear/vorticità rispetto a dati iniziali matched \cite{Raychaudhuri1955}. UMCH-U-0007: tali fonti non stabiliscono UMCH né le normalizzazioni del progetto.

\section{Norme e preregistrazione}\label{sec:u-norms}
Le sole norme preregistrate sono
\begin{equation}\label{eq:u-norms}\mathfrak C_2=\|\mathbf R\|_2,\qquad \mathfrak C_\infty=\|\mathbf R\|_\infty.\end{equation}
Norma, canale, regioni, finestre, sottrazioni e incertezze non possono cambiare dopo i dati. I canali possono essere correlati: le norme sono summary descrittive, non sei evidenze indipendenti e non definiscono una likelihood factorization. Senza un joint covariance/dependence model preregistrato lo stato è \verb|DEPENDENCE_UNRESOLVED| e le norme non possono confermare.

\section{Modello nullo e soglie}\label{sec:u-thresholds}
Ponendo $x=\ell/\ell_0\geq1$, confrontiamo sempre
\begin{align}\label{eq:u-families}
F_0&=0,&F_P&=Ax^{-p},&F_E&=Ae^{-q(x-1)},&F_{PE}&=A_\infty+A_1x^{-p}.
\end{align}
UMCH-U-0008: rumore o curvatura nota non sono prova di floor; ogni famiglia positiva compete con $F_0$.

\section{Controllo Minkowski}\label{sec:u-minkowski}
In Minkowski, con protocolli matched, $\mathbf R=0$. È il controllo nullo matematico e resta soluzione valida. La sua esistenza non confuta da sola una ipotesi limitata all'universo realizzato.

\section{Controllo FLRW}\label{sec:u-flrw}
Nel frame comoving FLRW, un controllo tidal è $R_{\rm tidal}=\ell^2\sqrt3|\ddot a/a|$. UMCH-U-0009: è test di definizione su curvatura cosmologica nota, non detection di un nuovo minimo.

\section{Controllo Schwarzschild}\label{sec:u-schwarzschild}
Nel frame statico dichiarato, Ricci è nullo ma gli autovalori tidal sono proporzionali a $(-2,1,1)$. Senza boundary data cosmologici il frame CMB nel vuoto resta unresolved; il frame statico del controllo non costituisce un frame universale.

\section{Controllo VSI}\label{sec:u-vsi}
Spaziotempi VSI possono avere invarianti scalari nulli con struttura di curvatura non banale \cite{PelavasEtAl2005}. UMCH-U-0010: un singolo invariante non è test sufficiente; il toy type-N valida il vettore operativo, non rappresenta dati.

\section{Teorema di non-identificabilità di $\ell_0$}\label{sec:u-no-go}
Per scale $\ell$ osservate,
\begin{align}\label{eq:u-degeneracy}
F_P&=(A\ell_0^p)\ell^{-p},\\
F_E&=(Ae^q)e^{-(q/\ell_0)\ell},\\
F_{PE}&=A_\infty+(A_1\ell_0^p)\ell^{-p}.
\end{align}
Quindi i dati identificano combinazioni, non $\ell_0$ separatamente; $F_0$ non contiene $\ell_0$. UMCH-U-0011: \verb|ELL0_STRUCTURALLY_NON_IDENTIFIABLE_UNDER_CURRENT_FAMILIES|. Più punti sulle stesse famiglie non eliminano una simmetria esatta di riparametrizzazione.

\section{Candidato con turnover}\label{sec:u-turnover}
Consideriamo, senza adottarlo come legge,
\begin{equation}\label{eq:u-turnover}F_T=A x^p e^{-q(x-1)}.
\end{equation}
Se $p,q$ sono fissati da teoria indipendente, due risposte positive danno
\begin{equation}\label{eq:u-inversion}\ell_0=\frac{q(\ell_1-\ell_2)}{p\log(\ell_1/\ell_2)-\log(y_1/y_2)},\end{equation}
e $\ell_{\rm peak}=(p/q)\ell_0$. UMCH-U-0012: è \verb|MATHEMATICAL_IDENTIFIABILITY_CANDIDATE_ONLY| e ramo \verb|SECONDARY_PROJECTION_SPECIFIC|; se $q$ è libero, la degenerazione ritorna. Il ramo operatoriale può in principio usare landmark proiettivi/spettrali senza q, ma nessuno è fisicamente derivato.

\section{Gate di derivazione della shape}\label{sec:u-shape}
La curvatura ha dimensione $L^{-2}$; per canali tidal, magnetic e holonomy area-normalizzati in regioni regolari piccole, $R\sim\ell^2|\mathrm{curvatura}|$, dunque $p=2$ nello scope. Per clock/null/congruence ciò richiede un protocollo short-baseline specifico. UMCH-U-0013: la GR locale non deriva $e^{-q\ell/\ell_0}$ né $q$. Il turnover è \verb|BLOCKED_PENDING_INDEPENDENT_NONLOCAL_MECHANISM|.

\section{Falsificazione e gate osservativo}\label{sec:u-falsification}
Una regione ammissibile, frame risolto e risposta compatibile con zero sotto incertezze preregistrate può contraddire una soglia positiva nello scope. Il gate ordinato richiede: assenza di \verb|FRAME_UNRESOLVED| e \verb|DOMAIN_INCONSISTENT|; famiglia prospectiva, altrimenti \verb|EXPLORATORY_FAMILY_SELECTION|; dipendenza, likelihood e nuisance risolti, altrimenti \verb|DEPENDENCE_UNRESOLVED|, \verb|LIKELIHOOD_UNRESOLVED| o \verb|NUISANCE_UNBOUNDED|; identificabilità e replica, altrimenti \verb|NON_IDENTIFIABLE| o \verb|REPLICATION_MISSING|. Anche tutti i pass producono solo \verb|CONFIRMATORY_ANALYSIS_ELIGIBLE_NOT_EVIDENCE|. UMCH-U-0014: eleggibilità non è evidenza e nessun requisito è sostituito da model fit favorevole.

\section{Ledger finale dei risultati}\label{sec:u-results}
\begin{center}\begin{tabular}{p{.48\linewidth}p{.43\linewidth}}\toprule
Oggetto&Stato\\\midrule
Core worldline&\verb|SUPERSEDED_AS_CORE|\\
Modelli fissi A/B&\verb|CONTRADICTED_UNDER_ASSUMPTIONS|\\
Core spacetime&\verb|UNPROVEN|\\
Frame vacuum senza certificato&\verb|FRAME_UNRESOLVED|\\
$\ell_0$ nelle famiglie correnti&\verb|ELL0_STRUCTURALLY_NON_IDENTIFIABLE_UNDER_CURRENT_FAMILIES|\\
$F_T$&\verb|MATHEMATICAL_IDENTIFIABILITY_CANDIDATE_ONLY|\\
Meccanismo per $q$&\verb|BLOCKED_PENDING_INDEPENDENT_NONLOCAL_MECHANISM|\\
Detection&\verb|NO_POSITIVE_DETECTION_CLAIM|\\\bottomrule\end{tabular}\end{center}

\section{Limiti e programma di ricerca}\label{sec:u-limits}
Restano i gate \verb|UMCH-A-0003|, \verb|UMCH-A-0004|, \verb|UMCH-A-0005|, \verb|UMCH-A-0008|, \verb|UMCH-A-0010|, \verb|UMCH-A-0011| e \verb|UMCH-A-0012|: frame domain-specific, normalizzazioni, likelihood/nuisance, shape, meccanismo e verifica umana. Topologia, caustiche e orizzonti richiedono scope dedicato. Il meccanismo nonlocale è \verb|EXTERNAL_EVIDENCE_REQUIRED| e la verifica indipendente è \verb|HUMAN_REVIEW_REQUIRED|. Il passo successivo non è scegliere un kernel conveniente: l'evoluzione autonoma si arresta finché questi input non sono forniti.

\section{Riproducibilità e assistenza IA}\label{sec:u-reproducibility}
Script deterministici riproducono casi esatti, soglie, gate frame, no-go e inversione turnover. Le versioni linguistiche condividono claim ID, equazioni e citazioni. L'IA ha assistito scrittura, codice e controllo; non è fonte né coautrice. Responsabilità e verifica scientifica restano umane.

\section{Conclusione}\label{sec:u-conclusion}
UMCH è stata corretta da vincolo pointwise worldline a programma operativo multiscala. Il contributo principale è delimitare rigorosamente ciò che è definito, falsificabile e non identificabile. UMCH-U-0015: non emerge evidenza positiva; emerge un percorso condizionale verso identificabilità, bloccato finché un meccanismo indipendente non fissa la shape.

\bibliographystyle{plain}\bibliography{../../../references/library}
\end{document}
